% This is a template for your written document.
%
% To compile using latexmk on the command line, run the following: 
% latexmk -pdf main.tex

\documentclass[12pt]{article}
\usepackage{setspace}
\usepackage{graphicx} % used for includegraphics
\singlespace
\usepackage[left=1in,right=1in,top=1in,bottom=1in]{geometry}

\title{\textbf{Design and Implementation of a Web-Based Stock Scoring Engine}}
\author{Leul Ejigu}

\begin{document}

\maketitle

\section*{Introduction}
Financial investment in the modern era is characterized by an huge  amount of data. Investors seeking to analyze the S\&P 500 index must navigate a complex landscape of quantitative metrics such as revenue growth and P/E ratios, alongside qualitative factors, such as leadership stability and competitive advantage. Visualizing data for over 500 companies presents a significant computational and design challenge. To address this, this project proposes the development of a web-based Stock Scoring Engine. This application will digest large financial datasets to produce a ranked list of stocks based on a proprietary 13-metrics scoring system(Thought it is essential for a stock to have ) .

The value of this Junior Is is in its approach to information hierarchy. Research by Dong et al. shows that visualizing financial "big data" works best when using specific "post-model" strategies. Instead of dumping raw and confusing data on the user, this system handles the heavy calculations in the backend. It processes the entire S\&P 500 and only presents the most important trends and summary stats for the top-ranked stocks. This keeps the dashboard clean and actionable, solving the information overload problem that usually makes financial analysis difficult for individual investors.

\begin{figure}[h]
\begin{center}
\includegraphics[scale=0.7]{methodology.png}
\caption{Archie}
\label{fig:method}       % Give a unique label
\end{center}
\end{figure}


\newpage
\section*{Appendix}
A concise list of features / user stories in the order in which they will be built. A few examples are below to demonstrate the expected scope and level of granularity; you will have more features than this.
\begin{itemize}
	\item Default picture display on web application.
	\item On a button-click, user can separate the image into foreground and background.
	\item User can select a picture from their desktop.
	\item Selected picture displays on the web application.
\end{itemize}


\bibliographystyle{acm}
\bibliography{bibliography.bib}

\end{document}
